\SourceLeaf{430}%
surface itself. For the radius of curvature we therefore have the simple formula
\[
\frac1r=Z(\xi\xi T+2\xi\eta U+\eta\eta V).
\]

\Needspace{5\baselineskip}
\begin{center}11.\end{center}
\noindent\hspace*{1.2em}Since infinitely many planes can be drawn through this normal, the radius of curvature may take infinitely many different values. In this connection $T,U,V,Z$ are to be regarded as constant and $\xi,\eta,\zeta$ as variable. To make the latter depend on a single variable, take two fixed points on the great circle whose pole is $L$, at a distance of $90^{0}$ from one another. Let them be called $M,M'$. Let their coordinates relative to the centre of the sphere be $\alpha,\beta,\gamma$; $\alpha',\beta',\gamma'$. We then have
\[
\cos\lambda(1)=\cos\lambda M.\cos M(1)+\cos\lambda M'.\cos M'(1)+\cos\lambda L.\cos L(1).
\]
Now if we set
\[
\lambda M=\varphi,
\]
then
\[
\cos\lambda M'=\sin\varphi
\]
and the formula becomes
\[
\xi=\alpha\cos\varphi+\alpha'\sin\varphi
\]
and likewise
\[
\begin{aligned}
\eta&=\beta\cos\varphi+\beta'\sin\varphi\\
\zeta&=\gamma\cos\varphi+\gamma'\sin\varphi.
\end{aligned}
\]

Thus if we set
\[
\begin{aligned}
A&=(\alpha\alpha T+2\alpha\beta U+\beta\beta V)Z\\
B&=(\alpha\alpha'T+(\alpha\beta'+\alpha\beta')U+\beta\beta'V)Z\\
C&=(\alpha'\alpha'T+2\alpha'\beta'U+\beta'\beta'V)Z
\end{aligned}
\]
we have:
